\begin{figure*}[b]
\centering
\begin{tikzpicture}
\draw[ge] (3.675,2.479) -- (3.170,3.412);
\draw[ge] (3.675,2.479) -- (4.378,3.297);
\draw[ge] (3.675,2.479) -- (2.428,2.570);
\draw[ge] (3.675,2.479) -- (4.207,1.350);
\draw[ge] (1.069,1.675) -- (2.109,1.572);
\draw[ge] (1.069,1.675) -- (2.163,0.522);
\draw[ge] (1.069,1.675) -- (2.428,2.570);
\draw[ge] (3.170,3.412) -- (4.378,3.297);
\draw[ge] (3.170,3.412) -- (2.428,2.570);
\draw[ge] (2.109,1.572) -- (2.428,2.570);
\draw[ge] (2.109,1.572) -- (3.156,1.536);
\draw[ge] (2.163,0.522) -- (3.156,1.536);
\draw[ge] (2.163,0.522) -- (3.424,0.522);
\draw[ge] (2.428,2.570) -- (3.156,1.536);
\draw[ge] (3.156,1.536) -- (3.424,0.522);
\draw[ge] (4.611,0.250) -- (4.207,1.350);
\draw[ge] (4.611,0.250) -- (3.424,0.522);
\draw[ge] (4.207,1.350) -- (3.424,0.522);
\draw[ge] (4.207,1.350) -- (4.917,2.213);
\node[gn, fill=pc0, text=white] (P1n1) at (3.675,2.479) {1};
\node[gn, fill=pc0, text=white] (P1n2) at (1.069,1.675) {2};
\node[gn, fill=pc0, text=white] (P1n3) at (3.170,3.412) {3};
\node[gn, fill=pc0, text=white] (P1n4) at (2.109,1.572) {4};
\node[gn, fill=pc1, text=white] (P1n5) at (2.163,0.522) {5};
\node[gn, fill=pc3, text=black] (P1n6) at (4.378,3.297) {6};
\node[gn, fill=pc0, text=white] (P1n7) at (2.428,2.570) {7};
\node[gn, fill=pc0, text=white] (P1n8) at (3.156,1.536) {8};
\node[gn, fill=pc1, text=white] (P1n9) at (4.611,0.250) {9};
\node[gn, fill=pc1, text=white] (P1n10) at (4.207,1.350) {10};
\node[gn, fill=pc1, text=white] (P1n11) at (3.424,0.522) {11};
\node[gn, fill=pc2, text=black] (P1n12) at (4.917,2.213) {12};
\node[ptitle] at (2.583,4.000) {(a) round 0 (Stage 1)};
\node[pstat] at (2.583,0.000) {4 comm., $Q{=}0.176$, NMI 0.509};
\node[barrow, rotate=0, minimum height=9mm] at (5.817,2.000) {};
\node[note, anchor=south] at (5.817,2.280) {round 1};
\draw[goodc, line width=1.1pt, rounded corners=2pt] (6.467,0.000) rectangle (11.633,4.000);
\draw[ge] (10.142,2.479) -- (9.637,3.412);
\draw[ge] (10.142,2.479) -- (10.845,3.297);
\draw[ge] (10.142,2.479) -- (8.894,2.570);
\draw[ge] (10.142,2.479) -- (10.673,1.350);
\draw[ge] (7.536,1.675) -- (8.576,1.572);
\draw[ge] (7.536,1.675) -- (8.630,0.522);
\draw[ge] (7.536,1.675) -- (8.894,2.570);
\draw[ge] (9.637,3.412) -- (10.845,3.297);
\draw[ge] (9.637,3.412) -- (8.894,2.570);
\draw[ge] (8.576,1.572) -- (8.894,2.570);
\draw[ge] (8.576,1.572) -- (9.623,1.536);
\draw[ge] (8.630,0.522) -- (9.623,1.536);
\draw[ge] (8.630,0.522) -- (9.891,0.522);
\draw[ge] (8.894,2.570) -- (9.623,1.536);
\draw[ge] (9.623,1.536) -- (9.891,0.522);
\draw[ge] (11.078,0.250) -- (10.673,1.350);
\draw[ge] (11.078,0.250) -- (9.891,0.522);
\draw[ge] (10.673,1.350) -- (9.891,0.522);
\draw[ge] (10.673,1.350) -- (11.383,2.213);
\node[gn, fill=pc0, text=white] (P2n1) at (10.142,2.479) {1};
\node[gn, fill=pc0, text=white] (P2n2) at (7.536,1.675) {2};
\node[gn, fill=pc0, text=white] (P2n3) at (9.637,3.412) {3};
\node[gn, fill=pc0, text=white] (P2n4) at (8.576,1.572) {4};
\node[gn, fill=pc0, text=white] (P2n5) at (8.630,0.522) {5};
\node[gn, fill=pc0, text=white] (P2n6) at (10.845,3.297) {6};
\node[gn, fill=pc0, text=white] (P2n7) at (8.894,2.570) {7};
\node[gn, fill=pc0, text=white] (P2n8) at (9.623,1.536) {8};
\node[gn, fill=pc1, text=white] (P2n9) at (11.078,0.250) {9};
\node[gn, fill=pc1, text=white] (P2n10) at (10.673,1.350) {10};
\node[gn, fill=pc1, text=white] (P2n11) at (9.891,0.522) {11};
\node[gn, fill=pc1, text=white] (P2n12) at (11.383,2.213) {12};
\node[ptitle] at (9.050,4.000) {(b) round 1: highest $Q$};
\node[pstat] at (9.050,0.000) {2 comm., $Q{=}0.253$, NMI 1.000};
\node[barrow, rotate=0, minimum height=9mm] at (12.283,2.000) {};
\node[note, anchor=south] at (12.283,2.280) {round 2};
\draw[ge] (16.608,2.479) -- (16.104,3.412);
\draw[ge] (16.608,2.479) -- (17.312,3.297);
\draw[ge] (16.608,2.479) -- (15.361,2.570);
\draw[ge] (16.608,2.479) -- (17.140,1.350);
\draw[ge] (14.002,1.675) -- (15.042,1.572);
\draw[ge] (14.002,1.675) -- (15.097,0.522);
\draw[ge] (14.002,1.675) -- (15.361,2.570);
\draw[ge] (16.104,3.412) -- (17.312,3.297);
\draw[ge] (16.104,3.412) -- (15.361,2.570);
\draw[ge] (15.042,1.572) -- (15.361,2.570);
\draw[ge] (15.042,1.572) -- (16.090,1.536);
\draw[ge] (15.097,0.522) -- (16.090,1.536);
\draw[ge] (15.097,0.522) -- (16.358,0.522);
\draw[ge] (15.361,2.570) -- (16.090,1.536);
\draw[ge] (16.090,1.536) -- (16.358,0.522);
\draw[ge] (17.545,0.250) -- (17.140,1.350);
\draw[ge] (17.545,0.250) -- (16.358,0.522);
\draw[ge] (17.140,1.350) -- (16.358,0.522);
\draw[ge] (17.140,1.350) -- (17.850,2.213);
\node[gn, fill=pc0, text=white] (P3n1) at (16.608,2.479) {1};
\node[gn, fill=pc0, text=white] (P3n2) at (14.002,1.675) {2};
\node[gn, fill=pc0, text=white] (P3n3) at (16.104,3.412) {3};
\node[gn, fill=pc0, text=white] (P3n4) at (15.042,1.572) {4};
\node[gn, fill=pc0, text=white] (P3n5) at (15.097,0.522) {5};
\node[gn, fill=pc0, text=white] (P3n6) at (17.312,3.297) {6};
\node[gn, fill=pc0, text=white] (P3n7) at (15.361,2.570) {7};
\node[gn, fill=pc0, text=white] (P3n8) at (16.090,1.536) {8};
\node[gn, fill=pc0, text=white] (P3n9) at (17.545,0.250) {9};
\node[gn, fill=pc0, text=white] (P3n10) at (17.140,1.350) {10};
\node[gn, fill=pc0, text=white] (P3n11) at (16.358,0.522) {11};
\node[gn, fill=pc0, text=white] (P3n12) at (17.850,2.213) {12};
\node[ptitle] at (15.517,4.000) {(c) round 2: TPOC-LPA output};
\node[pstat] at (15.517,0.000) {1 comm., $Q{=}0.000$, NMI 0.000};
\node[barrow, rotate=-90, minimum height=0.85cm] at (9.050,-1.0) {};
\node[note, anchor=west, font=\scriptsize\itshape] at (9.450,-1.0) {highest $Q$: passed to Stage~3};
\end{tikzpicture}
\caption{Best-iterate guard: Stage~2 rounds of TPOC-LPA on a planted graph (nodes 1--8 and 9--12). The guard passes the round with the highest $Q$ (green frame) to Stage~3.}
\label{fig:m2}
\end{figure*}
